\documentclass[12pt]{article}
\usepackage[T1]{fontenc}
\usepackage[utf8]{inputenc}
\usepackage{lmodern}
\usepackage{textcomp}
\usepackage{enumitem}
\usepackage{geometry}
\geometry{margin=1in}
\begin{document}
\begin{center}
Answer Key - Sociology - Graduate Level Assessment 3 \
\small (Answers are ordered exactly as in the question paper.)
\end{center}

\section*{Multiple Choice}
\begin{enumerate}[label=\arabic*.]
\item \textbf{Answer:} B
\\ \textit{Options:} A) charismatic authority; B) rational-legal authority; C) traditional authority; D) value-rational authority
\\ \textit{Note:} Correct option: B - rational-legal authority
\item \textbf{Answer:} A
\\ \textit{Options:} A) human behaviour is meaningful, and varies between individuals and cultures; B) it is difficult for sociologists to gain access to a research laboratory; C) sociologists are not rational or critical enough in their approach; D) we cannot collect empirical data about social life
\\ \textit{Note:} Correct option: A - human behaviour is meaningful, and varies between individuals and cultures
\item \textbf{Answer:} C
\\ \textit{Options:} A) secondary socialization and strict discipline; B) emotional support and sexual gratification; C) primary socialization and personality stabilization; D) oppressing women and reproducing the labour force
\\ \textit{Note:} Correct option: C - primary socialization and personality stabilization
\item \textbf{Answer:} A
\\ \textit{Options:} A) If all of your friends jumped off a bridge, I suppose you would too.; B) You should be proud to be a part of this organization.; C) If the door is closed, try the window.; D) Once a thief, always a thief.
\\ \textit{Note:} Correct option: A - If all of your friends jumped off a bridge, I suppose you would too.
\item \textbf{Answer:} B
\\ \textit{Options:} A) indoctrination and propaganda; B) freedom of movement for citizens; C) one-party rule; D) a centrally planned economy
\\ \textit{Note:} Correct option: B - freedom of movement for citizens
\end{enumerate}

\section*{True / False}
\begin{enumerate}[label=\arabic*.]
\setcounter{enumi}{5}
\item \textbf{Answer:} False
\\ \textit{Options:} A) True; B) False
\\ \textit{Note:} LLM-generated false statement. Correct answer: 'had begun to emerge through the separation of work and home life'.
\item \textbf{Answer:} True
\\ \textit{Options:} A) True; B) False
\\ \textit{Note:} LLM-generated true statement based on MCQ.
\item \textbf{Answer:} True
\\ \textit{Options:} A) True; B) False
\\ \textit{Note:} LLM-generated true statement based on MCQ.
\item \textbf{Answer:} True
\\ \textit{Options:} A) True; B) False
\\ \textit{Note:} LLM-generated true statement based on MCQ.
\item \textbf{Answer:} False
\\ \textit{Options:} A) True; B) False
\\ \textit{Note:} LLM-generated false statement. Correct answer: 'women's criminal behaviour tended to reflect traditional gender roles'.
\end{enumerate}

\section*{Long Form Response}
\begin{enumerate}[label=\arabic*.]
\setcounter{enumi}{10}
\item \textbf{Answer:} Functionalism. Explain why this is the correct answer and provide supporting evidence. Reference key facts or concepts that support this choice.
\\ \textit{Note:} Long-form derived from MCQ; award credit for stating the correct idea and giving supporting reasoning.
\item \textbf{Answer:} It specifies the poverty line at a level set in the 1960 s and adjusted since to reflect inflation.. Explain why this is the correct answer and provide supporting evidence. Reference key facts or concepts that support this choice.
\\ \textit{Note:} Long-form derived from MCQ; award credit for stating the correct idea and giving supporting reasoning.
\end{enumerate}

\vfill
\noindent\textit{End of Answer Key}
\end{document}